% !PRINTED 7
\weberpage{7}
%<<CONT
so besteht die Identität:
\weq{4}{a+b+c=0}
und wir erhalten die sechs Doppelverhältnisse
\weq{5}{\begin{aligned}
&-\frac{c}{b},\quad -\frac{a}{c},\quad -\frac{b}{a},\\
&-\frac{b}{c},\quad -\frac{c}{a},\quad -\frac{a}{b};
\end{aligned}}
setzen wir das erste von ihnen $-c:b=\varkappa^2$, so erhält man nach (4) die sechs:
\weq{6}{\varkappa^2,\quad 1-\frac{1}{\varkappa^2},\quad \frac{1}{1-\varkappa^2},\quad \frac{1}{\varkappa^2},\quad \frac{\varkappa^2}{\varkappa^2-1},\quad 1-\varkappa^2.}

Es sind also die sechs Doppelverhältnisse aus den vier Punkten \gesperrt{linear} durch eines unter ihnen ausgedrückt.

Wenn die $x_1, x_2, x_3, x_4$ untereinander permutiert werden, so werden die $a, b, c$ untereinander permutiert und ändern ihre Vorzeichen, jedoch so, daß entweder alle drei Vorzeichen gleichzeitig geändert werden oder ungeändert bleiben. Beispielsweise geben die Transpositionen

\[
\begin{aligned}
(14)\ \text{und}\ (23) & \begin{pmatrix} a & b & c \\ -a & -c & -b \end{pmatrix} \\
(24)\ \text{und}\ (31) & \begin{pmatrix} a & b & c \\ -c & -b & -a \end{pmatrix} \\
(34)\ \text{und}\ (12) & \begin{pmatrix} a & b & c \\ -b & -a & -c \end{pmatrix}
\end{aligned}
\]

% [sic] The comma after x_4 is printed; expected no comma before "reell".

Wenn $x_1, x_2, x_3, x_4$, reell sind, so wird $\varkappa^2$ dann und nur dann ein positiver echter Bruch, wenn $\varkappa^2$ und $1-\varkappa^2$ positiv, also entweder $a$ und $c$ positiv und $b$ negativ oder $a$ und $c$ negativ, $b$ positiv ist. Dies findet statt, wenn $x_1, x_2, x_3, x_4$ in dieser Reihenfolge der Größe nach aufsteigend einander folgen und bei allen den Anordnungen, die daraus durch solche Permutationen entstehen, die $b$ ungeändert lassen. Dies sind die folgenden acht:
\[
\begin{array}{cccc}
1 & 2 & 3 & 4 \\
2 & 3 & 4 & 1 \\
3 & 4 & 1 & 2 \\
4 & 1 & 2 & 3 \\
1 & 4 & 3 & 2 \\
4 & 3 & 2 & 1 \\
3 & 2 & 1 & 4 \\
2 & 1 & 4 & 3
\end{array}
\]
